\documentclass[11pt]{amsart}
\usepackage{amsmath,amssymb,booktabs,hyperref}
\hypersetup{pdfauthor={Otakhon U. Kenjaev},pdftitle={The two lost notes on page 54 of Ramanujan's lost notebook},pdfsubject={doi:10.5281/zenodo.23030496 (all versions)},pdfkeywords={Ramanujan, lost notebook, theta functions, cubic identities}}
\newtheorem{theorem}{Theorem}
\newtheorem{lemma}[theorem]{Lemma}
\theoremstyle{remark}\newtheorem{remark}[theorem]{Remark}
\newcommand{\f}{f}

\title{The two lost notes on page 54 of Ramanujan's lost notebook}
\author{Otakhon U. Kenjaev}
\date{September 29, 2026}

% STATUS v2 2026-09-29: Teorema 3-4 (circular summation + Borwein a). verifier 12/12, sonli 14/14. Son 2004 faqat zbMATH taqrizi.

\begin{document}
\begin{abstract}
On page 54 of his lost notebook Ramanujan states two families of cubic theta-function
identities and, in each family, writes ``see note'' next to one further claim
($A^3-B^3$ in the first family, $A^3+B^3$ in the second). Both notes are lost:
Andrews and Berndt record that they have not been preserved. We show that both
missing identities are special cases of another claim written on the \emph{same
page}, Ramanujan's cubic circular summation (Entry 8.2.3 in Andrews--Berndt), and
that in terms of the Borweins' cubic theta function
$a(q)=\sum_{m,n}q^{m^2+mn+n^2}$ they take the compact form
\[
A^3-B^3=q^2\f^3(-q^3,-q^{12})+\f(-q^2,-q^3)\,a(q^5),\qquad
q\,(A^3+B^3)=\f^3(-q^6,-q^9)-\f(-q,-q^4)\,a(q^5).
\]
We also record an elementary second form obtained from Ramanujan's other
identities on the page, and computational evidence that, unlike the four
recorded cube identities, the two missing combinations have no one- or two-term
theta-\emph{product} representation in a large explicit search space; the
appearance of the binary theta series $a(q^5)$, which is not such a product,
explains why these two claims were deferred to a note. All identities are
certified to $2000$ coefficients by a single integer-arithmetic script and
independently at 50-digit precision. (Version 2; version 1 contained only the
elementary form.)
\end{abstract}
\maketitle

\section{Introduction}
As usual, let
\[
\f(a,b):=\sum_{n=-\infty}^{\infty}a^{n(n+1)/2}\,b^{n(n-1)/2},\qquad |ab|<1,
\]
and $\f(-q):=\f(-q,-q^2)=(q;q)_\infty$. Page 54 of Ramanujan's lost notebook
\cite{LN} contains two groups of cubic identities for theta functions with
quotient of arguments $q^{15}$; they are treated as Entries 8.2.1 and 8.2.2 in
Andrews and Berndt's second volume \cite[pp.~176--178]{AB2}.

\textbf{Entry 8.2.1.} Let $A=\f(-q^7,-q^8)$ and $B=q\f(-q^2,-q^{13})$. Then
\begin{align}
A+B&=\frac{\f(-q^2,-q^3)}{\f(-q,-q^4)}\,\f(-q^5),\tag{8.2.6}\\
A-B&=\f(-q^{2/3},-q)+q^{2/3}\f(-q^3,-q^{12}),\tag{8.2.7}\\
A^3+B^3&=\frac{\f(-q^6,-q^9)}{\f(-q^3,-q^{12})}\,\f^3(-q^5),\tag{8.2.8}\\
(A-B)^3&=\frac{\f(-q^2,-q^3)\,\f^3(-q,-q^4)}{\f(-q^3,-q^{12})}+q^2\,\f^3(-q^3,-q^{12}).\tag{8.2.9}
\end{align}
After stating these, Ramanujan writes ``$A^3-B^3$ \emph{see note}.'' Andrews and
Berndt comment: ``Either Ramanujan never wrote this note, or if he did, it has
not been preserved'' \cite[p.~177]{AB2}.

\textbf{Entry 8.2.2.} Let $A=\f(-q^4,-q^{11})$ and $B=q\f(-q,-q^{14})$. Then
\begin{align}
A-B&=\frac{\f(-q,-q^4)}{\f(-q^2,-q^3)}\,\f(-q^5),\tag{8.2.10}\\
A+B&=-\frac{1}{q^{1/3}}\left\{\f(-q^{1/3},-q^{4/3})-\f(-q^6,-q^9)\right\},\tag{8.2.11}\\
A^3-B^3&=\frac{\f(-q^3,-q^{12})}{\f(-q^6,-q^9)}\,\f^3(-q^5),\tag{8.2.12}\\
(A+B)^3&=-\frac1q\left\{\frac{\f(-q,-q^4)\,\f^3(-q^2,-q^3)}{\f(-q^6,-q^9)}-\f^3(-q^6,-q^9)\right\}.\tag{8.2.13}
\end{align}
Here Ramanujan writes ``$A^3+B^3$ \emph{see note},'' and again ``he apparently
did not write this note, or, if he did, it has been lost'' \cite[p.~178]{AB2}.

To our knowledge the two missing notes have not been reconstructed in the
literature. The cubic identities of page 54 were first proved by Son
\cite{Son}, and Andrews and Berndt state that in this section they ``follow
precisely Son's beautiful proofs'' \cite[p.~173]{AB2}; both ``see note''
remarks, with the statements that the notes are lost, appear inside that very
section, and the only further content of \cite{Son} mentioned there is ``some
new modular equations'' \cite[p.~180]{AB2}. We were unable to find any later
source addressing the notes.

The main observation of this paper (Section~\ref{sec:circ}) is that the two
missing identities are hidden two lines further down the same page: they are
special cases of Ramanujan's cubic circular summation, Entry 8.2.3 of
\cite{AB2}, whose right-hand side is the Borweins' cubic theta function
$a(q)$. Section~\ref{sec:rekon} records a second, elementary form obtained from
(8.2.7), (8.2.9), (8.2.11), (8.2.13). Section~\ref{sec:yoqlik} shows that no
one- or two-term theta-product form exists within a large explicit search
space, and Section~\ref{sec:verifier} describes the verification scripts.

\section{The notes as circular summations}\label{sec:circ}
On page 54 Ramanujan also states the following, proved by Son \cite{Son04} and
given as Entry 8.2.3 in \cite[p.~178]{AB2}: for $|ab|<1$ and $q=ab$,
\begin{equation}\label{eq:circ}
\f^3(ab^2,a^2b)+a\f^3(b,a^3b^2)+b\f^3(a,a^2b^3)
=\f(a,b)\left\{\frac{\f^9(-q)}{\f^3(-q^3)}+27q\,\frac{\f^9(-q^3)}{\f^3(-q)}\right\}^{1/3}.
\end{equation}
With the Borweins' cubic theta functions \cite{BB91,BBG94}
\[
a(q)=\sum_{m,n=-\infty}^{\infty}q^{m^2+mn+n^2},\quad
b(q)=\frac{\f^3(-q)}{\f(-q^3)},\quad
c(q)=3q^{1/3}\frac{\f^3(-q^3)}{\f(-q)},\qquad a^3=b^3+c^3,
\]
the brace in \eqref{eq:circ} equals $(b^3(q)+c^3(q))^{1/3}=a(q)$.

\begin{theorem}\label{thm:3}
Let $A=\f(-q^7,-q^8)$ and $B=q\f(-q^2,-q^{13})$ (Entry 8.2.1). Then
\begin{equation}\label{eq:T3}
A^3-B^3=q^2\f^3(-q^3,-q^{12})+\f(-q^2,-q^3)\,a(q^5).
\end{equation}
\end{theorem}
\begin{proof}
Take $a=-q^3$, $b=-q^2$ in \eqref{eq:circ}, so $ab=q^5$. Then
$\f(ab^2,a^2b)=\f(-q^7,-q^8)=A$,
$a\f^3(b,a^3b^2)=-q^3\f^3(-q^2,-q^{13})=-B^3$,
$b\f^3(a,a^2b^3)=-q^2\f^3(-q^3,-q^{12})$ and $\f(a,b)=\f(-q^2,-q^3)$.
\end{proof}

\begin{theorem}\label{thm:4}
Let $A=\f(-q^4,-q^{11})$ and $B=q\f(-q,-q^{14})$ (Entry 8.2.2). Then
\begin{equation}\label{eq:T4}
A^3+B^3=\frac1q\left\{\f^3(-q^6,-q^9)-\f(-q,-q^4)\,a(q^5)\right\}.
\end{equation}
\end{theorem}
\begin{proof}
Take $a=-q^6$, $b=-q^{-1}$ in \eqref{eq:circ}, so again $ab=q^5$. Then
$\f(ab^2,a^2b)=\f(-q^4,-q^{11})=A$. Using $\f(a,b)=a\,\f(a^{-1},a^2b)$
\cite[p.~34]{BerndtIII}, $\f(-q^{-1},-q^{16})=-q^{-1}\f(-q,-q^{14})$ and
$\f(-q^{-1},-q^6)=-q^{-1}\f(-q,-q^4)$; hence
$a\f^3(b,a^3b^2)=-q^6\f^3(-q^{-1},-q^{16})=q^3\f^3(-q,-q^{14})=B^3$,
$b\f^3(a,a^2b^3)=-q^{-1}\f^3(-q^6,-q^9)$ and $\f(a,b)=-q^{-1}\f(-q,-q^4)$.
\end{proof}

Theorems~\ref{thm:3}--\ref{thm:4} have exactly the shape of the recorded
identities (8.2.9) and (8.2.13) --- a cube of a level-15 theta function plus one
further term --- except that the second term contains the binary theta series
$a(q^5)$ rather than a theta product. Since \eqref{eq:circ} is printed on the same
page, immediately after the two ``see note'' remarks, it is natural to read
Theorems~\ref{thm:3}--\ref{thm:4} as the content of the lost notes; this reading
is an interpretation, and we label it as such.

\emph{Literature.} We found no statement of Theorems~\ref{thm:3}--\ref{thm:4}, nor
any link between Entry 8.2.3 and the two ``see note'' remarks, in \cite{AB2}, in
the chapter on circular summation of \cite{AB3}, in \cite{AB4}, in
Chan--Liu--Ng \cite{CLN06} (full text), or in Ge--Luo \cite{GL19}. Son's paper
\cite{Son04} was checked only through its zbMATH review (Zbl~1062.11023), which
lists the general theorem and the cases $F_n$, $n=2,3,4,5,7$.

\section{An elementary second form}\label{sec:rekon}

\begin{theorem}\label{thm:1}
Let $A=\f(-q^7,-q^8)$ and $B=q\f(-q^2,-q^{13})$, as in Entry 8.2.1. Then
\begin{equation}\label{eq:T1}
A^3-B^3=\frac{\f(-q^2,-q^3)\,\f^3(-q,-q^4)}{\f(-q^3,-q^{12})}+q^2\,\f^3(-q^3,-q^{12})
+3AB\left\{\f(-q^{2/3},-q)+q^{2/3}\f(-q^3,-q^{12})\right\},
\end{equation}
where, by the Jacobi triple product,
\[
AB=q\,\f(-q^7,-q^8)\,\f(-q^2,-q^{13})
=q\,(q^2,q^7,q^8,q^{13};q^{15})_\infty\,(q^{15};q^{15})_\infty^2 .
\]
\end{theorem}

\begin{proof}
By the polynomial identity $A^3-B^3=(A-B)^3+3AB(A-B)$, the claim is precisely
(8.2.9) plus $3AB$ times (8.2.7), both of which are proved in
\cite[pp.~176--177]{AB2}. The product form of $AB$ is two applications of the
Jacobi triple product $\f(a,b)=(-a;ab)_\infty(-b;ab)_\infty(ab;ab)_\infty$.
\end{proof}

\begin{theorem}\label{thm:2}
Let $A=\f(-q^4,-q^{11})$ and $B=q\f(-q,-q^{14})$, as in Entry 8.2.2. Then
\begin{equation}\label{eq:T2}
A^3+B^3=-\frac1q\left\{\frac{\f(-q,-q^4)\,\f^3(-q^2,-q^3)}{\f(-q^6,-q^9)}-\f^3(-q^6,-q^9)\right\}
+\frac{3AB}{q^{1/3}}\left\{\f(-q^{1/3},-q^{4/3})-\f(-q^6,-q^9)\right\},
\end{equation}
where
\[
AB=q\,\f(-q^4,-q^{11})\,\f(-q,-q^{14})
=q\,(q,q^4,q^{11},q^{14};q^{15})_\infty\,(q^{15};q^{15})_\infty^2 .
\]
\end{theorem}

\begin{proof}
$A^3+B^3=(A+B)^3-3AB(A+B)$; substitute (8.2.13) and (8.2.11), proved in
\cite[p.~178]{AB2}.
\end{proof}

Both right-hand sides were verified coefficient-by-coefficient to $2000$ terms
(Section~\ref{sec:verifier}). We emphasize what is, and what is not, claimed:
the two theorems are elementary consequences of Ramanujan's own four recorded
identities, and in that sense they are certainly results Ramanujan possessed.
Comparing \eqref{eq:T1} with \eqref{eq:T3} gives, as a by-product,
$3AB(A-B)=\f(-q^2,-q^3)\{a(q^5)-\f^3(-q,-q^4)/\f(-q^3,-q^{12})\}$ for the pair of
Entry 8.2.1, and in the same way $3AB(A+B)=q^{-1}\f(-q,-q^4)\{a(q^5)-\f^3(-q^2,-q^3)/\f(-q^6,-q^9)\}$ for the pair of Entry 8.2.2.

\section{Why a note was necessary: a nonexistence computation}\label{sec:yoqlik}

Each of the four cube identities Ramanujan recorded on page 54 is a one- or
two-term combination of monomials in theta functions. It is natural to ask
whether $A^3-B^3$ (Entry 8.2.1) and $A^3+B^3$ (Entry 8.2.2) also possess such a
form, which Ramanujan would simply have recorded in the note. The evidence is
that they do not.

\subsection{Non-periodicity of the quadratic cofactors}
Write $A^3+B^3=(A+B)(A^2-AB+B^2)$ and $A^3-B^3=(A-B)(A^2+AB+B^2)$. A power
series with integer coefficients and constant term $1$ has a unique
representation $\prod_{n\ge1}(1-q^n)^{-c_n}$. If $(c_n)$ is periodic with
period $M$, the series equals $\prod_{a=1}^{M}(q^a;q^M)_\infty^{-c_a}$, a
generalized eta-quotient, i.e.\ a finite product of theta functions of the kind
occurring throughout page 54 (Jacobi triple product); conversely every such
product has periodic $(c_n)$. Computing $(c_n)$ for $n\le 380$
by the logarithmic-derivative recursion we find, in each entry, a sharp split:
\begin{center}
\begin{tabular}{lll}
\toprule
Entry & combination & exponent sequence $(c_n)$\\
\midrule
8.2.1 & $A^2-AB+B^2$ \ (factor of the recorded (8.2.8)) & periodic mod $15$\\
8.2.1 & $A^2+AB+B^2$ \ (factor of the lost $A^3-B^3$) & not periodic\\
8.2.2 & $A^2+AB+B^2$ \ (factor of the recorded (8.2.12)) & periodic mod $15$\\
8.2.2 & $A^2-AB+B^2$ \ (factor of the lost $A^3+B^3$) & not periodic\\
\bottomrule
\end{tabular}
\end{center}
(``Not periodic'' means: no period $\le 60$ holds for $n\le 380$; the periodic cases were checked from $n=1$.) Thus in each
family the recorded cube identity is the one whose quadratic cofactor is a
generalized eta-quotient, and the deferred one is the one whose cofactor is not. This
already rules out any identity of the shape
$A^3\mp B^3=(\text{linear factor known on p.~54})\times(\text{generalized eta-quotient})$.

\subsection{Systematic search over theta-monomial bases}
We searched for representations
\[
A^3-B^3=\sum_{i=1}^{r} c_i\, q^{s_i} M_i,\qquad r\le 2,\ c_i\in\{\pm1,\pm2,\pm3\},\ s_i\in\{0,\dots,3\},
\]
where each $M_i$ is a monomial of total degree $3$ with exponents in
$\{-1,1,2,3\}$ in at most three functions drawn from the $12$-function basis
\begin{gather*}
\f(-q,-q^4),\ \f(-q^2,-q^3),\ \f(-q^3,-q^{12}),\ \f(-q^6,-q^9),\ \f(-q^5),\\
\f(-q,-q^{14}),\ \f(-q^2,-q^{13}),\ \f(-q^4,-q^{11}),\ \f(-q^7,-q^8),\ \f(-q),\ \f(-q^3),\ \f(-q^{15})
\end{gather*}
--- every theta function occurring in Entries 8.2.1--8.2.2 together with the
relevant $\eta$-type functions ($2344$ monomials; hash-matching on the first
$80$ coefficients, surviving candidates confirmed on $260$). The search was
repeated in the $q^{1/3}$-lattice (setting $q\mapsto q^3$ and adjoining the
functions $\f(-q^{a/3},-q^{b/3})$ occurring in (8.2.7) and (8.2.11): $20$ basis
functions, $11800$ monomials, shifts $s_i\in\{0,\dots,9\}$), and with the
three-term shape $c_1 q^{s_1}M_1+c_2 q^{s_2}M_2\pm3q^{s}AB\,M_3$ suggested by
Theorems~\ref{thm:1}--\ref{thm:2}.

\emph{Control of the instrument.} Run on the recorded combinations, the same
search re-finds Ramanujan's identities exactly: for $(A-B)^3$ of Entry 8.2.1 it
returns the two-term representation (8.2.9), and in the $q^{1/3}$-lattice it
returns (8.2.7) for $A-B$.

\emph{Result.} For $A^3-B^3$ (Entry 8.2.1) and $A^3+B^3$ (Entry 8.2.2) the
search returns nothing beyond trivial rewritings of the definitions. This is
consistent with Theorems~\ref{thm:3}--\ref{thm:4}: the missing identities need
the binary theta series $a(q^5)$, which lies outside the product basis. Within the
declared space --- which contains every identity actually appearing on page 54
--- the two lost combinations have no one- or two-term product form.

We state the scope honestly: this is a finite search over an explicit basis and
coefficient range, not a proof of nonexistence in general; a completeness
argument via the theory of modular functions on $\Gamma_0(15)$ (the natural
next step) is left open.

\section{The mirror symmetry, and a remark on the note}\label{sec:mirror}
The two entries are mirror images: under the exchange
$\f(-q^3,-q^{12})\leftrightarrow\f(-q^6,-q^9)$,
$\f(-q,-q^4)\leftrightarrow\f(-q^2,-q^3)$, identity (8.2.8) corresponds to
(8.2.12) and (8.2.9) to (8.2.13). In each family Ramanujan recorded the cube
combination with the generalized eta-quotient cofactor and deferred the other to a note.
The parallel structure of the two deferrals --- the \emph{same} obstruction on
both sides of the mirror --- is explained by Theorems~\ref{thm:3}--\ref{thm:4}:
both notes are the same circular summation \eqref{eq:circ}, with the same
modulus $ab=q^5$ and the same factor $a(q^5)$, applied to the two mirror pairs. Whether
Ramanujan intended precisely this, or a representation of a shape we have not
searched, cannot of course be decided; we label this paragraph as
interpretation, not theorem.

\section{The verification script}\label{sec:verifier}
All twelve identities of this paper --- (8.2.6)--(8.2.13),
\eqref{eq:T1}--\eqref{eq:T2} and \eqref{eq:T3}--\eqref{eq:T4} --- are certified by one self-contained Python
script \texttt{verifier.py} ($152$ lines, no dependencies). It works in
the variable $y=q^{1/3}$ with exact integer arithmetic; every identity
involving a quotient is verified after cross-multiplication, so no division of
power series occurs. Each identity is checked to $400$ coefficients of $q$ by
default, and the run reported here used $2000$ coefficients ($6003$ powers of
$y$; total runtime about one second). The script prints one MOS/FARQ
(match/mismatch) line per identity and exits nonzero on any mismatch; the
archived run prints $12/12$ MOS. A second, independent script \texttt{sonli\_tekshir.py} evaluates all identities as printed (with division and $q^{1/3}$) at five values of $q$ with 50-digit arithmetic; the largest relative discrepancy is $2\cdot10^{-49}$. A third script \texttt{ildiz.py} checks the foundations independently: it compares every theta series used with its Jacobi triple product (3000 coefficients), verifies the circular summation \eqref{eq:circ} itself as an identity of two-variable integer power series in independent variables $a,b$ (all monomials $a^ib^j$ with $i+j\le150$, where the brace is replaced by $a(ab)$), and checks $a^3=b^3+c^3$ in the division-free form $a(q)^3(q;q)_\infty^3(q^3;q^3)_\infty^3=(q;q)_\infty^{12}+27q(q^3;q^3)_\infty^{12}$ (3000 coefficients). Finally, \texttt{ildiz\_gp.gp} re-implements the theta series, Theorems~\ref{thm:3}--\ref{thm:4}, both by-products and the Borwein identity from scratch in PARI/GP (1500 coefficients), independently of the Python code.

\section*{Data availability}
The verification script and the search scripts of Section~\ref{sec:yoqlik} are
archived at Zenodo: \url{https://doi.org/10.5281/zenodo.23030496} (all versions; version~1 is doi:10.5281/zenodo.23030497).

\begin{thebibliography}{9}
\bibitem{LN} S.~Ramanujan, \emph{The Lost Notebook and Other Unpublished Papers}, Narosa, New Delhi, 1988.
\bibitem{AB2} G.~E.~Andrews and B.~C.~Berndt, \emph{Ramanujan's Lost Notebook, Part II}, Springer, New York, 2009 (Chapter 8).
\bibitem{BerndtIII} B.~C.~Berndt, \emph{Ramanujan's Notebooks, Part III}, Springer, New York, 1991 (Chapter 20, Entries 10(i)--(iv)).
\bibitem{Son04} S.~H.~Son, \emph{Circular summations of theta functions in Ramanujan's lost notebook}, Ramanujan J. \textbf{8} (2004), 235--272, doi:10.1023/B:RAMA.0000040483.55191.d1.
\bibitem{AB3} G.~E.~Andrews and B.~C.~Berndt, \emph{Ramanujan's Lost Notebook, Part III}, Springer, New York, 2012.
\bibitem{AB4} G.~E.~Andrews and B.~C.~Berndt, \emph{Ramanujan's Lost Notebook, Part IV}, Springer, New York, 2013.
\bibitem{BB91} J.~M.~Borwein and P.~B.~Borwein, \emph{A cubic counterpart of Jacobi's identity and the AGM}, Trans. Amer. Math. Soc. \textbf{323} (1991), 691--701, doi:10.1090/S0002-9947-1991-1010408-0.
\bibitem{BBG94} J.~M.~Borwein, P.~B.~Borwein and F.~G.~Garvan, \emph{Some cubic modular identities of Ramanujan}, Trans. Amer. Math. Soc. \textbf{343} (1994), 35--47, doi:10.1090/S0002-9947-1994-1243610-6.
\bibitem{CLN06} H.~H.~Chan, Z.-G.~Liu and S.~T.~Ng, \emph{Circular summation of theta functions in Ramanujan's lost notebook}, J. Math. Anal. Appl. \textbf{316} (2006), 628--641, doi:10.1016/j.jmaa.2005.05.015.
\bibitem{GL19} J.-K.~Ge and Q.-M.~Luo, \emph{Some extensions for Ramanujan's circular summation formula}, arXiv:1901.09820 (2019).
\bibitem{Son} S.~H.~Son, \emph{Cubic identities of theta functions}, Ramanujan J. \textbf{2} (1998), 303--316, doi:10.1023/A:1009751614537.
\end{thebibliography}
\end{document}
